\section{Human annotated dataset details}\label{appendix:dataset_details}
We use the human annotated data collected by~\textcite{August2024KnowYA}. The dataset includes 60 summaries over 10 papers, 6 summaries per paper. Of the 6 summaries, 2 are written by experts and 4 are machine written by GPT3. The 10 papers were sampled from the top 10\% of papers from \textit{r/science}, a subreddit dedicated to public discussions of scientific papers. These papers were chosen as a proxy for scientific topics the general public is most interested in. The dataset was annotated by 593 Mechanical Turk workers in total across the three tasks in the original study.
\Cref{tab:human_annotation_distribution} contains the distributions of scores assigned by the human annotators.

\input{research-chapters/5-chapter/tables/human_annotation_distribution}

In order to measure inter-annotator agreement, we bin the scores into a binary ``high-readability'' and ``low readability.'' Summaries given scores of 3 or higher are considered highly readable while summaries assigned scores less than 3 are considered to have low readability. We use Cohen's Kappa to calculate an inter-annotator agreement of 0.6. This is a moderate agreement for a somewhat subjective task, indicating that there is some general notion of readability. We also note that this is significantly higher than the agreement between traditional metrics and LMs (0.17 as shown in~\Cref{subsec:results-rqfour}).



\section{LM readability evaluation prompts}\label{appendix:prompt}
We experiment with 3 prompts, shown in~\Cref{tab:prompts}. The \texttt{Simple Prompt} simply asks the LM to rate the text for reading ease on a scale of 1 to 5. The American Society for Cell Biology (ASCB) provides guidelines for best practices in scientific communication.\footnote{\href{https://www.ascb.org/science-policy-public-outreach/science-outreach/communication-toolkits/best-practices-in-effective-science-communication/}{ASCB Best Practices in Science Communication
}} In the  \texttt{ASCB Prompt}, we provide these guidelines to the LM as guidance for rating the readability. Finally, the  \texttt{Own Reasoning Prompt} is similar to the  \texttt{Simple Prompt}, but with the additional instruction for the LM to use its own judgment to rate the text, rather than traditional readability formulas, such as FKGL. 

\begin{table}[!t]
\centering
    \begin{subtable}{.49\textwidth}
        \centering \footnotesize
       \begin{tabular}{r|lll}
        Model & Simple & ASCB & Own \\ \hline
        Mistral 7B & 0.46 & 0.54 & 0.52 \\
        Mixtral 7B & 0.46 & 0.47 & 0.54 \\
        Gemma 1.1 7B & 0.55 & 0.33 & 0.54 \\
        Llama 3.1 8B & 0.54 & 0.56 & 0.45 \\
        Llama 3.3 70B & 0.59 & 0.58 & 0.56 \\ \hdashfull
        Mean Corr. & 0.52 & 0.50 & 0.52
        \end{tabular}
        \vspace{-1mm}
        \caption{Pearson Correlation.}
        \label{subtab:prompt-comparison-pearson}
    \end{subtable}    
    \hfill
    \begin{subtable}{.49\textwidth}
        \centering \footnotesize
           \begin{tabular}{r|lll}
            Model & Simple & ASCB & Own \\ \hline
            Mistral 7B & 0.32 & 0.40 & 0.44 \\
            Mixtral 7B & 0.36 & 0.41 & 0.41 \\
            Gemma 1.1 7B & 0.42 & 0.24 & 0.43 \\
            Llama 3.1 8B & 0.38 & 0.35 & 0.34 \\
            Llama 3.3 70B & 0.36 & 0.38 & 0.35 \\ \hdashfull
            Mean Corr. & 0.37 & 0.36 & 0.39
            \end{tabular}
        \vspace{-1mm}
        \caption{Kendall-Tau Correlation.}\label{subtab:prompt-comparison-kt}
    \end{subtable} 
    \caption[Correlation with human judgment for each prompt.]{Pearson and Kendall-Tau Correlation with human judgment for each prompt listed in \Cref{tab:prompts}. \texttt{Own Reasoning} prompt performs the best averaged across all models.}
    \label{tab:prompt-comparison}
\end{table}


We report the Pearson and Kendall-Tau correlation of each prompt with human judgment in \Cref{tab:prompt-comparison}. The \texttt{Own Reasoning Prompt} performs the best when averaged across all models. We found that the models tended to over-rely on the guidance provided in the \texttt{ASCB Prompt}, providing lower scores if the conditions are not met. For the \texttt{Simple Prompt}, the models would occasionally try to calculate FKGL or another readability metric, rather than using their own reasoning. This is likely because FKGL is strongly associated with readability in the models' training data. We found that the \texttt{Own Reasoning Prompt} struck the right balance between providing enough instructions that the model is able to understand the task without providing too much information for the model to over-rely on. However, it is notable that the \texttt{ASCB Prompt}, the worst performing prompt, still achieves higher correlation with human judgment than FKGL, the most popular traditional metric. 



\begin{table*}[]
\centering \scriptsize
\begin{tabular}{l}
% \hline
\textbf{Simple Prompt} \\ \hline
\begin{tabular}[c]{p{.9\textwidth}}
\vspace{-1mm} 
On a scale of  1 to 5, what is the reading ease of the following text? 1 indicates the text requires expert background knowledge and 5 indicates the text is readable to the general population. Assume the reader is an adult.\textnewline \textnewline \\ Format the output as follows:\textnewline\\ Score: \textless{}score\textgreater\textnewline\\ Reason: \textless{}reasoning\textgreater \textnewline \\ Text: \{SUMMARY\}
\vspace{5mm}
\end{tabular} \\

% \hline
\textbf{ASCB Guidelines Prompt} \\ \hline
\begin{tabular}[c]{p{.9\textwidth}}
\vspace{-1mm} 
On a scale of  1 to 5, what is the reading ease of the following text? 1 indicates the text requires expert background knowledge and 5 indicates the text is readable to the general population. Characteristics of a highly readable text include:\textnewline\\ - Know your audience, and focus and organize your information for that particular audience.\textnewline\\ - Focus on the big picture. What larger problem is your work a part of? What major ideas or issues does your work address? How will your work help global understanding of some issue?\textnewline\\ - Avoid jargon. If you must use a technical term, make sure to explain it, but simplify the language.\textnewline\\ - Try to use metaphors or analogies to everyday experiences that people can relate to.\textnewline\\ - Underscore the importance of public support for exploratory research and scientific information, and the role of this information in providing the context for effective policy making.\textnewline\textnewline\\ Assume the reader is an adult. Do not use Flesch-Kincaid or other readability formulas. Use your own judgment to rate the text.\textnewline\textnewline\\ Format the output as follows:\textnewline\\ Score: \textless{}score\textgreater\textnewline\\ Reason: \textless{}reasoning\textgreater\textnewline\textnewline\\ Text: \{SUMMARY\}
\vspace{5mm}
\end{tabular} \\

% \hline
\textbf{Own Reasoning Prompt} \\ \hline
\begin{tabular}[c]{p{.9\textwidth}}
\vspace{-1mm}
On a scale of  1 to 5, what is the reading ease of the following text? 1 indicates the text requires expert background knowledge and 5 indicates the text is readable to the general population. \textnewline\\ Assume the reader is an adult. Do not use Flesch-Kincaid or other readability formulas. Use your own judgment to rate the text.\textnewline\textnewline\\ Format the output as follows:\textnewline\\ Score: \textless{}score\textgreater\textnewline\\ Reason: \textless{}reasoning\textgreater \textnewline\textnewline\\ Text: \{SUMMARY\}
% \vspace{5mm}
\end{tabular} \\
% \hline
\end{tabular}
\caption[Prompts we tested.]{Prompts we tested. \texttt{Own Reasoning} is the best performing prompt, as reported in~\Cref{tab:prompt-comparison}.}\label{tab:prompts}
\end{table*}
